\documentclass[preprint,12pt,numbers]{elsarticle}
\usepackage{amsmath,amssymb,amsthm}
\usepackage[hidelinks]{hyperref}
\newtheorem{theorem}{Theorem}[section]
\newtheorem{lemma}[theorem]{Lemma}
\newtheorem{proposition}[theorem]{Proposition}
\newtheorem{corollary}[theorem]{Corollary}
\makeatletter
\let\ps@pprintTitle\ps@plain
\makeatother
\begin{document}
\begin{frontmatter}
\title{Blocking vertical translates of cyclotomic Frobenius graphs}
\begin{abstract}
Let $q=s^n$ and let $m\ge2$ divide $s-1$. Partition the multiplicative
group of $\mathbb F_q$ into $m$ cosets. On each coset multiply the
Frobenius map $x\mapsto x^s$ by a nonzero element of $\mathbb F_s$.
We determine the order of the minimum number of vertical translates of
the resulting graph needed to meet every affine line. If the ratios of
the branch coefficients generate a group of order $h$, and at least two
coefficients occur, the minimum is of order $ns$ when $h$ divides $n$,
and of order $\log q$ otherwise. The comparison constants depend only
on $m$. The lower bounds come from blocking trace hyperplanes and from
simultaneous multiplicative character prescriptions with a growing
number of translates. Two nonmonomial permutations on the same field
of order $t^6$ can have blocking numbers of orders $t^2$ and $\log q$,
even though each uses only two cosets.
\end{abstract}
\begin{keyword}
finite field \sep blocking set \sep cyclotomic mapping \sep Frobenius map \sep multiplicative character sum
\MSC[2020] 11T06 \sep 11T23 \sep 51E21
\end{keyword}
\end{frontmatter}

\section{The blocking number and the coefficient period}
Let $f:E\to E$ be a function on a finite field. For $S\subseteq E$ consider
the graphs $y=f(x)-\alpha$, $\alpha\in S$, together with the point
$[0:1:0]$ at infinity. Write $B(f)$ for the least size of a nonempty
$S$ for which this union meets every $E$-rational projective line.
Vertical lines and the line at infinity contain the distinguished point.
For a line $y=ax+b$ the condition is equivalent to
\begin{equation}\label{eq:block}
   (b+S)\cap (f-a\,\mathrm{id})(E)\ne\varnothing
   \qquad(a,b\in E).
\end{equation}
Equivalently, the same set of parameters must cover the intercepts
at every slope. Asgarli, Ghioca and Yip studied this question for
$f(x)=x^d$ with $d$ fixed and $|E|$ growing \cite[Section~5]{AGY}.
Here we let the degree vary even among permutations of a single field.

Let $F=\mathbb F_s\subset E=\mathbb F_q$, $q=s^n$, and fix $m\ge2$ with
$m\mid s-1$. The map $\lambda(x)=x^{(q-1)/m}$ sends $E^\times$ onto
the group $\mu_m\subseteq F^\times$ of $m$th roots of unity. Given a table
$\beta:\mu_m\to F^\times$, put
\begin{equation}\label{eq:graph}
  f(0)=0,\qquad f(x)=\beta(\lambda(x))x^s\quad(x\ne0).
\end{equation}
Write $C=\operatorname{im}\beta$ and $r=|C|$. For any $c_0\in C$ the subgroup
generated by $c/c_0$, $c\in C$, is generated by all the quotients $c/d$,
$c,d\in C$. Denote its order by $h$. The familiar cyclotomic
description of a permutation \cite{BW} is particularly simple here:
since $\mu_m\subseteq F$, the Frobenius map preserves each label, and
the coset of label $z$ is mapped bijectively to that of
$z\lambda(\beta(z))$. Thus $f$ is a permutation precisely when the
induced map on labels is bijective, an instance of Wang's cyclotomic
criterion \cite{Wang}. The bounds below apply whenever
at least two coefficients occur, whether or not $f$ is a permutation.

\begin{theorem}\label{thm:main}
Suppose $r\ge2$. Uniformly in $s,n$ and the coefficient table,
\[
 B(f)=\begin{cases}
         \Theta_m(ns),&h\mid n,\\
         \Theta_m(\log q),&h\nmid n .
       \end{cases}
\]
The index $m$ alone controls the implied constants. More explicitly,
with $K_m=2m(1+(1+2m^2)^2)+1$, the first case satisfies
\[
  \frac{1+n(s-1)}{r}\le B(f)\le K_m ns.
\]
In the second case $B(f)\le\lceil2m\log(2q)\rceil$. For either case,
if $\log q\ge16384m^5$, then
\[
                 B(f)>\frac{\log q}{32m^2\log m}.
\]
All logarithms are natural.
\end{theorem}

For example, the monomial $x^{s+(q-1)/m}$ is obtained by taking
$\beta(z)=z$. Then $h=m$. For a general table $h$ can exceed the
number of branches, even when $m=2$. The difference comes from a
nonzero slope at which every branch becomes a singular linear map.
Section~\ref{sec:positive} characterizes all such slopes and proves the upper
bounds. Section~\ref{sec:escape} gives a missing line for any sufficiently
small set of translates. Together these prove Theorem~\ref{thm:main}:
for the bounded range below the displayed threshold one may use
$B(f)\ge1$ and decrease the lower comparison constant. When $s$ is fixed,
$ns$ is itself of order $\log q$. The separation of growth rates is
visible as $s$ grows, as in Section~\ref{sec:examples}.

\section{Trace hyperplanes and random affine lines}
\label{sec:positive}

Write $N=N_{E/F}$, $\operatorname{Tr}=\operatorname{Tr}_{E/F}$ and
$\sigma(x)=x^s$. For $c\in C$ and $a\in E$ put $L_{c,a}(x)=c\sigma(x)-ax$.
A nonzero element of its kernel satisfies $x^{s-1}=a/c$. In the cyclic
group $E^\times$ the $(s-1)$st powers form $\ker N$, so $L_{c,a}$ is singular
precisely when $N(a)=c^n$. In this case its kernel is $wF$ for any nonzero
kernel vector $w$. For a fixed $c_0\in C$, the equalities $c^n=c_0^n$
for all $c\in C$ are equivalent to $h\mid n$, since the subgroup generated
by the ratios $c/c_0$ is cyclic of order $h$.

\begin{proposition}\label{prop:trace-hyperplanes}
If $h\mid n$, then $rB(f)\ge n(s-1)+1$.
\end{proposition}
\begin{proof}
Choose $v_c\in E^\times$ with $v_c^{s-1}=c_0/c$. This is possible since
$N(c_0/c)=1$. Put $w_c=(c_0v_c)^{-1}$. These $r$ vectors are linearly
independent over $F$: they are eigenvectors of $\sigma$ with distinct
eigenvalues $c/c_0$. For $t\in E^\times$ choose the slope $a=c_0t^{s-1}$.
Then
\[
 L_{c,a}(tv_cz)=c_0t^sv_c(\sigma(z)-z),\qquad
 \operatorname{im} L_{c,a}
  =\{y\in E:\operatorname{Tr}(t^{-s}w_c y)=0\}.
\]
Indeed, $\sigma-1$ has a one-dimensional kernel and its image lies in
the trace-zero hyperplane. The trace pairing is nondegenerate since $E/F$ is separable.

Let $S$ be a blocking parameter set. If each projection
$\operatorname{Tr}(t^{-s}w_c S)$ omitted an element $z_c\in F$, independence
of the normals and nondegeneracy of the trace pairing would give
$b\in E$ with $\operatorname{Tr}(t^{-s}w_cb)=-z_c$ for every $c$.
Then $b+S$ would be disjoint from every $\operatorname{im} L_{c,a}$,
contradicting blocking at the line $y=ax+b$. Consequently
$\bigcup_{c\in C}w_c S$ projects onto $F$ under every map
$y\mapsto\operatorname{Tr}(t^{-s}y)$. As $t$ varies these are all the
nonzero linear forms on $E$ over $F$. This union therefore meets every
affine hyperplane of $F^n$. Its cardinality is at least $n(s-1)+1$ by
the affine blocking theorem of Jamison and Brouwer--Schrijver
\cite{BS,Jamison}, and at most $r|S|$.
\end{proof}

\begin{proposition}\label{prop:separated-upper}
If $h\nmid n$, then $B(f)\le\lceil 2m\log(2q)\rceil$.
\end{proposition}
\begin{proof}
For every $a\in E$, at least one $c\in C$ satisfies $c^n\ne N(a)$.
The corresponding $L_{c,a}$ is injective. Restrict it to a nonzero
multiplicative class carrying coefficient $c$, together with zero.
This gives $(q-1)/m+1\ge q/m$ distinct values of $f(x)-ax$.
A uniform parameter thus hits each fixed line with probability at least
$1/m$. For $k=\lceil 2m\log(2q)\rceil$ independent parameters the expected
number of missed slope--intercept pairs is at most $q^2e^{-k/m}\le1/4$.
Some choice hits all pairs.
\end{proof}

\begin{lemma}\label{lem:line-sampler}
Let $H$ be a subgroup of index $m$ in $E^\times$, and let
$L(x)=x^s-ax$. Choose $Y\in E$ and $V\in E^\times$ independently and
uniformly. For every $b\in E$,
\[
 \Pr\bigl((Y+FV)\cap (L(H)-b)=\varnothing\bigr)
 \le \frac{1+2m^2}{s}.
\]
\end{lemma}
\begin{proof}
Translation invariance allows $b=0$. For $A\subseteq E$ of density $\rho$,
put $X_A=|(Y+FV)\cap A|$. Its mean is $\mu=s\rho$.
For distinct $i,j\in F$, the map $(Y,V)\mapsto(Y+iV,Y+jV)$
is a bijection onto the ordered pairs of distinct points of $E$.
Thus the two membership indicators have covariance
$-\rho(1-\rho)/(q-1)$, and
\[
 \operatorname{Var}X_A=s\rho(1-\rho)\frac{q-s}{q-1}\le\mu .
\]
If $L$ is invertible, use $A=L(H)$, of density
$\rho=(q-1)/(mq)\ge1/(2m)$. Chebyshev's inequality gives
$\Pr(X_A=0)\le1/\mu$.

Suppose $L$ is singular, with kernel $K=wF$ and image the hyperplane
$U$. Apply the variance formula to $A=H$. Conditional on any $V=v$,
the mean over $Y$ is still $\mu$. Multiplication by $u\in H$ carries
$Y+Fv$ to $uY+Fuv$ and preserves $H$, so the conditional variance is
constant for $v$ in each of the $m$ multiplicative cosets of $H$.
Those cosets have equal mass, and the average conditional variance is
at most $\mu$. In particular,
\[
 \Pr_Y\bigl((Y+K)\cap H=\varnothing\bigr)
 \le \frac{m}{\mu}.
\]
The left-hand side is exactly the fraction of $U$ missed by $L(H)$:
$L(Y)$ is uniform on $U$ and its fibre is $Y+K$.
The sampled direction lies in $U$ with probability
$(q/s-1)/(q-1)<1/s$. Conditional on any other direction $v$, the sampled
line meets $U$ in one uniform point, since translation by every element
of $U$ preserves its distribution. The total probability of missing
$L(H)$ is consequently at most $1/s+m/(s\rho)\le(1+2m^2)/s$.
\end{proof}

\begin{proposition}\label{prop:uniform-upper}
Write $A_m=1+2m^2$. Every coefficient table in the theorem satisfies
\[
 B(f)\le \bigl(2m(1+A_m^2)+1\bigr)ns .
\]
If $s\ge A_m^2$, the sharper bound $(4n+1)s$ holds.
\end{proposition}
\begin{proof}
Scale the ordinate by $\beta(1)^{-1}$, which preserves the minimum by
scaling slopes, intercepts and vertical parameters together. The resulting
function agrees with $x\mapsto x^s$ on $H=\ker\lambda$ and at zero.
For a fixed slope $a$ and intercept $b$, a set of shifts meeting
$L(H)-b$ hits the corresponding line. By Lemma~\ref{lem:line-sampler},
one random affine $F$-line of shifts fails with probability at most
$A_m/s$. When $s\ge A_m^2$, the expected number of pairs missed by
$4n+1$ independent such lines is at most
$q^2s^{-(4n+1)/2}=s^{-1/2}<1$. Their union has at most $(4n+1)s$ points.

When $s<A_m^2$, the map $L$ has fibres of size at most $s$, including
the singular case. Its domain $H\cup\{0\}$ has at least $q/m$ elements.
Every fixed line is therefore hit by a uniform shift with probability at
least $1/(ms)$. The point sampling of Proposition~\ref{prop:separated-upper}
gives $B(f)\le\lceil 2ms\log(2q)\rceil$. Since
$\log(2q)=\log 2+n\log s\le n(1+A_m^2)$, this gives the asserted
bound; the same bound also dominates $(4n+1)s$ for $m\ge2$.
\end{proof}

\section{A logarithmic lower bound from simultaneous omissions}\label{sec:escape}

For each $c\in C$ write $J_c=\beta^{-1}(c)$ for its set of labels.

\begin{proposition}\label{prop:escape}
Suppose $r\ge2$. If $\log q\ge16384m^5$, every $S\subset E$ of
cardinality $\ell\le\log q/(32m^2\log m)$ admits a missing line.
The coefficients, including their multiplicities in the table, and the
extension degree may vary with $q$.
\end{proposition}

In particular the logarithmic lower bound holds for both sides of the
divisibility criterion. We first record the counting estimate used in
the three ranges of $n$.

\begin{lemma}\label{lem:pattern}
Let $K=\mathbb F_Q$, let $\chi$ be a complex multiplicative character of
exact order $m$, and let $P_1,\ldots,P_M\in K[T]$ be nonzero. Denote by
$D$ the number of distinct roots of their product over $\overline K$.
Suppose that, for each nonzero $e\in\{0,\ldots,m-1\}^M$,
$\prod_iP_i^{e_i}$ is not an $m$th power in $\overline K[T]$.
For prescribed $m$th roots of unity $\xi_i$, let $A$ count those $t\in K$
for which all $P_i(t)$ are nonzero and $\chi(P_i(t))=\xi_i$ for every $i$.
Then
\[
 \left|A-\frac{Q}{m^M}\right|\le D(\sqrt Q+1).
\]
\end{lemma}

\begin{proof}
The assertion for $M=0$ is immediate. Delete the $z\le D$ rational zeros
of the product. On the remaining points use
$m^{-1}\sum_{j=0}^{m-1}(\chi(P_i(t))/\xi_i)^j$ for the indicator of
the $i$th prescription. The zero exponent vector contributes
$(Q-z)/m^M$. Each other complete sum has absolute value at most
$(D-1)\sqrt Q$ by the multiplicative Weil bound
\cite[Theorem~5.41]{LN}, applied after division by the leading
coefficient. Restriction to the chosen points costs at most $z$ in that
sum (extend $\chi$ to zero by zero). The $m^M-1$ coefficients of these
sums have absolute value $m^{-M}$ each. Together with the principal
term this gives the stated estimate.
\end{proof}

\begin{proof}[Proof of Proposition~\ref{prop:escape}]
We may suppose $S\ne\varnothing$. Identify the cyclic value group of
$\lambda$ with the complex $m$th roots of unity, using one fixed
isomorphism. Prescriptions on $\lambda$ can then be counted by
Lemma~\ref{lem:pattern}.

\paragraph{Large extension degree.}
Suppose $n\ge4m$. Put $\sigma(x)=x^s$ and $L_c(x)=c\sigma(x)-ax$.
For $a\ne0$, $L_c$ is singular on $E$ exactly when
$N_{E/F}(a)=c^n$: its nonzero kernel equation is
$x^{s-1}=a/c$, and the image of this power map is the norm-one group.
In the singular case the kernel has $s$ elements. If
$\{c^n:c\in C\}\ne F^\times$, choose the norm of $a$ outside that set.
Otherwise $s-1\le r\le m\le s-1$, so $s=m+1$, and the values $c^n$
are all distinct. Any $a\ne0$ then gives exactly one singular branch.

Let $T(x)=a^{-1}x^s$, now also on $\overline E$. The $F$-linear
operators $U_c=cT-1$ commute. Use products to denote composition and put
\[
 P=\prod_{c\in C}U_c,\qquad Q_c=\prod_{d\in C\setminus\{c\}}U_d .
\]
For a prospective parameter $t\in E$ set $b=aP(t)$. For each
$\alpha\in S$, $c\in C$, and every $u\in E$ satisfying $L_c(u)=\alpha$,
form the polynomial
\[
                 H_{\alpha,c,u}(T_0)=Q_c(T_0)+u .
\]
Indeed $L_c(Q_c(t))=aP(t)$, so the entire rational fibre above
$b+\alpha$ consists of these candidates on branch $c$.
Choose one prescribed label outside $J_c$ for each candidate. Such a
label exists since $r\ge2$. If the prescriptions hold, each candidate
is nonzero and its actual label selects a coefficient different from
$c$. Hence none contributes to $f(x)-ax=b+\alpha$. It is essential here
to index the product by the distinct coefficients and to exclude the
whole set $J_c$. An invertible branch supplies one candidate per shift.
The possible singular branch supplies zero or $s=m+1$ candidates,
according as its fibre is empty or nonempty. Thus $M\le2m\ell$ even in
this exceptional case. Each $H_{\alpha,c,u}$ has degree $s^{r-1}$ and
derivative $(-1)^{r-1}$.

We check the geometric nonpower condition, including the case of
composite $m$. Every zero of $H_{\alpha,c,u}$ lies in $aP(t)=-\alpha$,
so candidates attached to different shifts have disjoint root sets.
Choose a point $t_0$ of this geometric fibre; $P$ is nonconstant.
The eigenspace $V_c=\ker U_c$ is a line over $F$: a nonzero eigenvector
satisfies $v^{s-1}=a/c$, and the ratios of two such vectors lie in
$F^\times$. The spaces $V_c$ have distinct $T$-eigenvalues $c^{-1}$.
Their direct sum has $s^r$ elements, all in $\ker P$. The polynomial
$P$ has degree $s^r$, hence the geometric fibre consists exactly of
the points $t_0+\sum_c v_c$ with $v_c\in V_c$.

On $V_c$, $Q_c$ acts as multiplication by
$w_c=\prod_{d\ne c}(d/c-1)\ne0$; on each $V_d$, $d\ne c$, it vanishes.
As $L_c(Q_c(t_0)+u)=0$, the zero set of $H_{\alpha,c,u}$ on the fibre
has the single coordinate equation
$v_c=-w_c^{-1}(Q_c(t_0)+u)\in V_c$. Distinct $u$ on this branch specify
distinct values. If a product of the candidates with exponents between
$0$ and $m-1$ is an $m$th power, its root multiplicity at every one of
these points vanishes modulo $m$. This says
$\sum_c h_c(v_c)=0$ in $\mathbb Z/m\mathbb Z$, where $h_c$ at a
prescribed coordinate is its exponent, and at every other coordinate is
zero. Varying one coordinate shows that $h_c$ is constant. Every
invertible branch prescribes only one of its $s$ coordinates, so its
constant is zero. The sole remaining constant, if any, vanishes by the
same sum, including when all $s$ coordinates were prescribed. Doing
this in each shift fibre kills every exponent. Lemma~\ref{lem:pattern}
therefore applies over $E$ with $D\le Ms^{r-1}$.

\paragraph{Bounded extension degrees greater than one.}
Suppose $2\le n<4m$. We justify below the inequalities
$s-1>m+m^2$ and $(nr\ell)^2<s$. Choose
\[
 N_{E/F}(a)\notin\{c^n:c\in C\}\cup\{cd:c,d\in C\}.
\]
All $L_c$ are invertible. Write $A_c=L_c^{-1}$, and take $b=u+t$ with
$u\in E$ and $t\in F$. The unique candidate on branch $c$ for a shift
$\alpha$ becomes $A_c(u+\alpha)+tA_c(1)$. Its norm is represented by
the polynomial
\[
 G_{c,\alpha}(Z)=
 \prod_{j=0}^{n-1}\bigl(\sigma^j A_c(u+\alpha)+Z\sigma^j A_c(1)\bigr)
 \in F[Z].
\]
Frobenius permutes the factors. Since $A_c(1)\ne0$, the polynomial
has degree $n$. For $x\in E^\times$ we have
\[
  \lambda(x)=N_{E/F}(x)^{(s-1)/m}.
\]
Thus a character of order $m$ on $F$ detects the prescribed label.

To separate the roots, first note that no two operators
$\sigma^j A_c$, indexed by $(j,c)$, are $E$-proportional. After
conjugating, a proportionality would read $A_c=v\sigma^k A_d$ for
$v\ne0$, $0\le k<n$. Applying $L_c$ and substituting $L_d(y)$ gives
\[
 d\sigma(y)-ay
   =c\sigma(v)\sigma^{k+1}(y)-av\sigma^k(y).
\]
The maps $1,\sigma,\ldots,\sigma^{n-1}$ are $E$-linearly independent:
a nonzero linearized polynomial of degree at most $s^{n-1}$ cannot
vanish on $s^n$ points. If $k=0$, comparison gives $v=1$, $c=d$.
Otherwise the two exponent supports can agree only for $n=2$, $k=1$;
then $d=-av$ and $-a=c\sigma(v)$ force $a\sigma(a)=cd$, which we
excluded.

Write $U=\sigma^j A_c$. Its normalized map $U/U(1)$ is distinct from
all the other normalized maps and is invertible over $F$. The roots
in question are $-U(u+\alpha)/U(1)$. For the same $U$ and different
shifts they never coincide. For different $U$ a coincidence restricts
$u$ to an empty set or a coset of a proper $F$-linear kernel, of size
at most $q/s$. Since $(nr\ell)^2<s$, there exists $u$ for which all the
roots are distinct. Each $G_{c,\alpha}$ thus has $n$ simple roots
shared by no other member. At any such root the multiplicity of a
Fourier product is the exponent of that member. It follows that no
nonzero exponent vector yields a geometric $m$th power.
Lemma~\ref{lem:pattern} now applies over $F$ with $M=r\ell$ and
$D=nM$.

\paragraph{The ground field.}
If $n=1$, choose $a\in C$. Each label with coefficient $a$ maps to zero
under $f-a\,\mathrm{id}$; each of the other labels maps onto one
multiplicative coset of $\ker\lambda$. At most $m-1$ nonzero cosets
occur. Specify one missing coset for all $t+\alpha$, $\alpha\in S$.
These $\ell$ linear polynomials have distinct simple roots, so
Lemma~\ref{lem:pattern} applies over $F=E$ with $D=M=\ell$.

\paragraph{Uniformity as the number of translates grows.}
Let $Q$ be the field size in the corresponding application of the lemma
and set $X=\log Q$. In the large-degree range $Q=q$ and
$M\log m\le X/(16m)$. In the other ranges $Q=s$, $n<4m$ and
$M\log m\le m\ell\log m\le nX/(32m)\le X/8$.
In every case $X\ge4096m^4$, so $M\le X/(8\log2)$.
For bounded $n$, the numerical conditions used above follow from
\[
 s-1\ge X>m+m^2,\qquad
 nr\ell\le 4m^2 X,\qquad
 (4m^2 X)^2<X^3/6<e^X=s.
\]
Here $\ell\le X$ follows from its stipulated range and $n<4m$.
Also $n\le4m\le 1+X/4\le Q^{1/4}$. For large $n$ the degree of a
candidate satisfies $s^{r-1}\le q^{1/4}$. Thus $D\le M Q^{1/4}$
in every range, and the error in Lemma~\ref{lem:pattern} is at most
$2M Q^{3/4}$. The main term is at least $Q^{7/8}$. It exceeds the
error since $2M\le X/(4\log2)<X/2<e^{X/8}$ for $X\ge256$.
Consequently there is an admissible $t$. Its associated intercept gives
the required missing line.
\end{proof}

\section{Two branches and examples on one field}\label{sec:examples}
Let $h\ge3$ and suppose $s\equiv1\pmod{2h}$. In odd characteristic take
$\gamma\in F^\times$ of order $h$. On the nonzero squares take the
coefficient $1$ and on the nonsquares the coefficient $\gamma$, so that
\[
 f(x)=x^s\left(\frac{1+\gamma}{2}
              +\frac{1-\gamma}{2}x^{(q-1)/2}\right).
\]
Since $\gamma$ is a square in $F$, both branches preserve their labels.
This gives a permutation whose coefficient period is $h$. It cannot
be a scalar multiple of a monomial function: on $E^\times$ the ratios
$f(x)/x^s$ would then form a coset of a multiplicative subgroup.
The two-element set $\{1,\gamma\}$ can do so only if $\gamma=-1$.

For any prime power $t\equiv1\pmod 6$, take $q=t^6$ and choose
$\gamma$ of order three in $\mathbb F_t$. With $s=t^2$, the relative
degree is three and Theorem~\ref{thm:main} gives $B(f)=\Theta(t^2)$.
With $s=t^3$, the relative degree is two and $B(f)=\Theta(\log q)$.
The coefficients and the two multiplicative cosets have stayed the same.
Since both functions permute $E$, every horizontal line already meets
either graph.

\begin{corollary}
For odd $s$, let $q=s^2$ and $f(x)=x^{s+(q-1)/2}$. Then $B(f)=s$.
Every one-dimensional $F$-subspace of $E$ is an optimal parameter set.
\end{corollary}
\begin{proof}
Here $r=h=n=2$. Proposition~\ref{prop:trace-hyperplanes} gives
$2B(f)\ge 2s-1$. Fix any $F$-line $W$ through zero. All elements of
$F^\times$ are squares in $E$, so $V_a=(f-a\,\mathrm{id})(E)$ is
closed under $F^\times$-scaling. Each square class spans $E$ over $F$,
since it has $(s^2-1)/2>s-1$ elements. Hence the $F$-span of $V_a$
contains the images of $\sigma-a$ and $-\sigma-a$. Their difference
is $2\sigma$, which is onto. The image of $V_a$ in $E/W$ thus contains
a nonzero element, all its $F^\times$-multiples, and zero. This proves
$V_a-W=E$ for every slope.
\end{proof}

For $n=1$ one can sharpen the lower coefficient. At a slope equal to
one branch coefficient, choose a missing multiplicative coset as in
Section~\ref{sec:escape}. If a set of $\ell$ shifts blocks the graph,
no intercept lies in that coset after all $\ell$ translations.
Lemma~\ref{lem:pattern} applied to the distinct linear polynomials
$T+\alpha$ gives $q/m^\ell\le\ell(\sqrt q+1)$.
In particular $\liminf B(f)/\log q\ge 1/(2\log m)$ for fixed $m$ in
this stratum. Indeed, if $\ell\le\log q$, taking logarithms gives
$\ell\log m\ge\tfrac12\log q-\log\log q-\log(1+q^{-1/2})$.
If $\ell>\log q$ the stated lower coefficient follows from
$2\log m>1$. Determining the best constant for arbitrary relative
degree remains open.

\section*{Verification and data availability}
OpenAI's Codex assisted in deriving the proofs and translating their
finite-field and counting arguments into Lean. The project records
the kernel audit in \texttt{lean/verification.txt}, including the
hashes of the sources and the axioms used by each exported theorem.
The accompanying Lean~4.8 development constructs an actual minimizing
set of translates and verifies both inequalities in the case $h\mid n$.
For $h\nmid n$ the sole external mathematical input in its minimum
theorem is the classical one-variable multiplicative Weil estimate,
applied to characters of exact order $m$ over the field of summation.
The finite preimages, geometric nonpower conditions, root counts, and
numerical estimates for growing $\ell$ are included in the development.
The formal counting argument allows an error $W+2z$, where $W$ bounds
each nontrivial complete sum and $z$ counts excluded rational inputs.
For $n=1$ this gives $\sqrt q+2$ in place of $\sqrt q+1$ in the finite
bound above. The formal sampling upper bound is $2m\log(2q)+1$.
The proofs and independent finite-field checks can be
reproduced from the
\href{https://github.com/Heinz1729/monomial-blocking-curves}{public repository}.
The script \texttt{checks/branch\_extensions.py} verifies all 62,500
prescribed fibre comparisons in one four-label, three-coefficient
example over $\mathbb F_{125}$. On $\mathbb F_{7^6}$ it checks all
117,648 nonzero slopes in each of the two order-three examples. It finds
2,451 simultaneous degeneracies for $s=49$ and zero for $s=343$.

\section*{Declaration of generative AI and AI-assisted technologies}
OpenAI's Codex was used to develop and check proofs, write Lean
formalizations and finite computations, and draft and edit this text.
The repository records the AI-assisted development under the name
GPT-Astra. The mathematical inputs to the verification are specified
above. This public version precedes approval by a responsible human author.

\begin{thebibliography}{9}
\raggedright
\bibitem{AGY} S. Asgarli, D. Ghioca, C.H. Yip, Blocking sets from a union of plane curves,
Des. Codes Cryptogr. 94 (2026), article 217.
\href{https://doi.org/10.1007/s10623-026-01958-z}{doi:10.1007/s10623-026-01958-z}.
\bibitem{BW} A. Bors, Q. Wang, Generalized cyclotomic mappings: switching between polynomial,
cyclotomic, and wreath product form, Commun. Math. Res. 38 (2022), 246--318.
\href{https://doi.org/10.4208/cmr.2021-0029}{doi:10.4208/cmr.2021-0029}.
\bibitem{Wang} Q. Wang, Cyclotomic mapping permutation polynomials over finite fields,
in: Sequences, Subsequences, and Consequences, Lecture Notes in Computer Science
4893, Springer, Berlin, 2007, 119--128.
\href{https://doi.org/10.1007/978-3-540-77404-4_11}{doi:10.1007/978-3-540-77404-4\_11}.
\bibitem{BS} A.E. Brouwer, A. Schrijver, The blocking number of an affine space,
J. Combin. Theory Ser. A 24 (1978), 251--253.
\href{https://doi.org/10.1016/0097-3165(78)90013-4}{doi:10.1016/0097-3165(78)90013-4}.
\bibitem{Jamison} R.E. Jamison, Covering finite fields with cosets of subspaces,
J. Combin. Theory Ser. A 22 (1977), 253--266.
\href{https://doi.org/10.1016/0097-3165(77)90001-2}{doi:10.1016/0097-3165(77)90001-2}.
\bibitem{LN} R. Lidl, H. Niederreiter, Finite Fields, 2nd ed.,
Cambridge University Press, Cambridge, 1997.
\href{https://doi.org/10.1017/CBO9780511525926}{doi:10.1017/CBO9780511525926}.
\end{thebibliography}
\end{document}
